\section{Descent of independent parameters}
\label{sec:descent}

The main theorem adjoins to the base a tuple $e$ independent from the configuration.
First-order consequences of the theorem can be brought back to the original base $k$
by the following lemma, which is the universal-domain form of the ``place''
argument of Evans--Hrushovski (1991, Lemma 2.1.1) and of the coheir argument used by
Bays--Zou.

\begin{lemma}[density of $k$-points]\label{lem:kpoints-dense}
Let $U$ be an irreducible $k$-variety and $f\in K[X]\setminus I_U$.  Then there is
$u_0\in U(k)=U(K)\cap k^N$ with $f(u_0)\ne0$.
\end{lemma}
\begin{proof}
Choose a $k$-basis $(\beta_j)$ of $K$ and write $f=\sum_j f_j\beta_j$ with $f_j\in k[X]$;
since $I_U$ is defined over $k$, $f\notin I_U$ means some $f_j\notin I_U\cap k[X]=:J$,
a prime ideal of $k[X]$.  If $f_j$ vanished on $U(k)=Z_k(J)$ then, by the Nullstellensatz
over the algebraically closed field $k$ (\lean{MvPolynomial.vanishingIdeal\_zeroLocus\_eq\_radical}
with both fields equal to $k$), $f_j\in\sqrt J=J$, a contradiction.  So some
$u_0\in U(k)$ has $f_j(u_0)\neq0$; as the $\beta_j$ are linearly independent over $k$,
$f(u_0)\ne0$.
\end{proof}

\begin{theorem}[descent]\label{thm:descent}
Let $\Phi\subseteq K^{N}\times K^{M}$ be constructible over $k$, let $w\in K^M$ and let
$e\in K^N$ with $e\ind{}w$ and $\Phi(e,w)$.  Then there is $e_0\in k^N$ with $\Phi(e_0,w)$.
\end{theorem}
\begin{proof}
Let $U=\Loc(e)$ and $C=\{u\in U(K):\Phi(u,w)\}$, a set constructible over $k(w)$
containing $e$.  Since $e\ind{}w$, $e$ is a generic point of $U$ over $\acl(w)$
(\cref{lem:dim}(2)).  By \cref{lem:gen-constructible} (with $L=k(w)$), $C$ contains
every generic point of $U$ over $\acl(w)$, so its complement $U(K)\setminus C$, which is
constructible over $k(w)$ and contains no generic point, is contained in $Z(f)$ for some
$f\in k(w)[X]\setminus I_U$.  By \cref{lem:kpoints-dense} there is $u_0\in U(k)$ with
$f(u_0)\ne0$; then $u_0\in C$, i.e.\ $\Phi(u_0,w)$.
\end{proof}

\begin{remark}[how to use descent; its limits]\label{rem:descent-limits}
Any consequence of \cref{thm:main} written as ``$\Phi(e,abcxyz)$'' for a
constructible $\Phi$ over $k$ (with $e$ the adjoined parameters, and all further
ambient-field witnesses existentially quantified) therefore holds with parameters
in $k$. To use this for interalgebraicity or interdefinability, fix the actual
finite polynomial or quasi-rational witness formulas supplied by the proof,
including their coefficients in $k$ and their nonvanishing conditions. A fixed
polynomial equation and a fixed polynomial inequation define constructible sets;
finite Boolean combinations preserve constructibility, as does existential
quantification over finite tuples of $K$ when $K$ is algebraically closed. For
example, a chosen relation
$P(e,v,u)=0$, together with the nonvanishing of a chosen coefficient of $P$ as a
polynomial in $u$, gives a constructible algebraicity witness. A chosen relation
$Q(e,v)u^{p^d}=R(e,v)$ with $Q(e,v)\ne0$ gives a constructible definability witness
in positive characteristic; use exponent $1$ in characteristic zero.

A degree bound alone does not make the full relations ``$u\in\acl(k,e,v)$'' or
``$u\in\dcl(k,e,v)$'' constructible in the ambient tuples. Already with empty
$(e,v)$, the roots of nonzero polynomials over $k$ of degree at most one form
exactly $k$. Constant rational-function witnesses also give exactly $k$. This
proper infinite subfield is not a constructible subset of the line: constructible
subsets are finite or cofinite, while the complement of $k$ contains $t+k$ for
any $t\notin k$. Quantifying coefficients over $K$ does not restrict them to $k$.
The earlier wording ``once a degree bound is fixed'' is refuted in issue \#31;
its original provenance is \texttt{6b76e01}, this remark. The Lean statements
\lean{GroupConfig.not\_isConstructibleOver\_range\_algebraMap} and
\lean{GroupConfig.not\_isConstructibleOver\_setOf\_exists\_natDegree\_le\_one}
record the base-field obstruction and its degree-one algebraicity case. The
corrected fixed-formula statement is given by
\lean{GroupConfig.isConstructibleOver\_iff\_exists\_polynomial\_data},
\lean{GroupConfig.IsConstructibleOver.polynomial\_zero},
\lean{GroupConfig.IsConstructibleOver.polynomial\_ne\_zero},
\lean{GroupConfig.QuasiRational.isConstructibleOver\_graph}, the finite Boolean
operations and \lean{GroupConfig.IsConstructibleOver.exists}.

The assertions ``$\rk(u/e)=r$'', ``$G$ is connected'' and ``$\rk(G)=n$'' are not
first-order in the parameter as stated; definability of Morley rank and degree
in $\mathrm{ACF}$ treats families, but we do not prove that here. The application
in \cref{sec:application} still requires fixing formulas for the
interalgebraicity and interdefinability conclusions to be descended, and proving
that the needed independence of the descended witnesses follows from the rank
equalities of the original configuration. These obligations are not discharged
by a degree bound or by the constructible-relation descent theorem alone.
\end{remark}
